% Options for packages loaded elsewhere
% Options for packages loaded elsewhere
\PassOptionsToPackage{unicode}{hyperref}
\PassOptionsToPackage{hyphens}{url}
\PassOptionsToPackage{dvipsnames,svgnames,x11names}{xcolor}
%
\documentclass[
  12pt,
  letterpaper,
  DIV=11,
  numbers=noendperiod]{scrartcl}
\usepackage{xcolor}
\usepackage[margin=1in]{geometry}
\usepackage{amsmath,amssymb}
\setcounter{secnumdepth}{5}
\usepackage{iftex}
\ifPDFTeX
  \usepackage[T1]{fontenc}
  \usepackage[utf8]{inputenc}
  \usepackage{textcomp} % provide euro and other symbols
\else % if luatex or xetex
  \usepackage{unicode-math} % this also loads fontspec
  \defaultfontfeatures{Scale=MatchLowercase}
  \defaultfontfeatures[\rmfamily]{Ligatures=TeX,Scale=1}
\fi
\usepackage{lmodern}
\ifPDFTeX\else
  % xetex/luatex font selection
  \setmainfont[]{Work Sans}
  \setmonofont[]{Arial}
\fi
% Use upquote if available, for straight quotes in verbatim environments
\IfFileExists{upquote.sty}{\usepackage{upquote}}{}
\IfFileExists{microtype.sty}{% use microtype if available
  \usepackage[]{microtype}
  \UseMicrotypeSet[protrusion]{basicmath} % disable protrusion for tt fonts
}{}
\makeatletter
\@ifundefined{KOMAClassName}{% if non-KOMA class
  \IfFileExists{parskip.sty}{%
    \usepackage{parskip}
  }{% else
    \setlength{\parindent}{0pt}
    \setlength{\parskip}{6pt plus 2pt minus 1pt}}
}{% if KOMA class
  \KOMAoptions{parskip=half}}
\makeatother
% Make \paragraph and \subparagraph free-standing
\makeatletter
\ifx\paragraph\undefined\else
  \let\oldparagraph\paragraph
  \renewcommand{\paragraph}{
    \@ifstar
      \xxxParagraphStar
      \xxxParagraphNoStar
  }
  \newcommand{\xxxParagraphStar}[1]{\oldparagraph*{#1}\mbox{}}
  \newcommand{\xxxParagraphNoStar}[1]{\oldparagraph{#1}\mbox{}}
\fi
\ifx\subparagraph\undefined\else
  \let\oldsubparagraph\subparagraph
  \renewcommand{\subparagraph}{
    \@ifstar
      \xxxSubParagraphStar
      \xxxSubParagraphNoStar
  }
  \newcommand{\xxxSubParagraphStar}[1]{\oldsubparagraph*{#1}\mbox{}}
  \newcommand{\xxxSubParagraphNoStar}[1]{\oldsubparagraph{#1}\mbox{}}
\fi
\makeatother


\usepackage{longtable,booktabs,array}
\newcounter{none} % for unnumbered tables
\usepackage{calc} % for calculating minipage widths
% Correct order of tables after \paragraph or \subparagraph
\usepackage{etoolbox}
\makeatletter
\patchcmd\longtable{\par}{\if@noskipsec\mbox{}\fi\par}{}{}
\makeatother
% Allow footnotes in longtable head/foot
\IfFileExists{footnotehyper.sty}{\usepackage{footnotehyper}}{\usepackage{footnote}}
\makesavenoteenv{longtable}
\usepackage{graphicx}
\makeatletter
\newsavebox\pandoc@box
\newcommand*\pandocbounded[1]{% scales image to fit in text height/width
  \sbox\pandoc@box{#1}%
  \Gscale@div\@tempa{\textheight}{\dimexpr\ht\pandoc@box+\dp\pandoc@box\relax}%
  \Gscale@div\@tempb{\linewidth}{\wd\pandoc@box}%
  \ifdim\@tempb\p@<\@tempa\p@\let\@tempa\@tempb\fi% select the smaller of both
  \ifdim\@tempa\p@<\p@\scalebox{\@tempa}{\usebox\pandoc@box}%
  \else\usebox{\pandoc@box}%
  \fi%
}
% Set default figure placement to htbp
\def\fps@figure{htbp}
\makeatother





\setlength{\emergencystretch}{3em} % prevent overfull lines

\providecommand{\tightlist}{%
  \setlength{\itemsep}{0pt}\setlength{\parskip}{0pt}}



 


\KOMAoption{captions}{tableheading}
\makeatletter
\@ifpackageloaded{caption}{}{\usepackage{caption}}
\AtBeginDocument{%
\ifdefined\contentsname
  \renewcommand*\contentsname{Table of contents}
\else
  \newcommand\contentsname{Table of contents}
\fi
\ifdefined\listfigurename
  \renewcommand*\listfigurename{List of Figures}
\else
  \newcommand\listfigurename{List of Figures}
\fi
\ifdefined\listtablename
  \renewcommand*\listtablename{List of Tables}
\else
  \newcommand\listtablename{List of Tables}
\fi
\ifdefined\figurename
  \renewcommand*\figurename{Figure}
\else
  \newcommand\figurename{Figure}
\fi
\ifdefined\tablename
  \renewcommand*\tablename{Table}
\else
  \newcommand\tablename{Table}
\fi
}
\@ifpackageloaded{float}{}{\usepackage{float}}
\floatstyle{ruled}
\@ifundefined{c@chapter}{\newfloat{codelisting}{h}{lop}}{\newfloat{codelisting}{h}{lop}[chapter]}
\floatname{codelisting}{Listing}
\newcommand*\listoflistings{\listof{codelisting}{List of Listings}}
\makeatother
\makeatletter
\makeatother
\makeatletter
\@ifpackageloaded{caption}{}{\usepackage{caption}}
\@ifpackageloaded{subcaption}{}{\usepackage{subcaption}}
\makeatother
\usepackage{bookmark}
\IfFileExists{xurl.sty}{\usepackage{xurl}}{} % add URL line breaks if available
\urlstyle{same}
\hypersetup{
  pdftitle={Iowa Trauma Registry Anomaly Detection Report},
  pdfauthor={Iowa Department of Health and Human Services},
  colorlinks=true,
  linkcolor={blue},
  filecolor={Maroon},
  citecolor={Blue},
  urlcolor={Blue},
  pdfcreator={LaTeX via pandoc}}


\title{Iowa Trauma Registry Anomaly Detection Report}
\usepackage{etoolbox}
\makeatletter
\providecommand{\subtitle}[1]{% add subtitle to \maketitle
  \apptocmd{\@title}{\par {\large #1 \par}}{}{}
}
\makeatother
\subtitle{Annual Case Reporting Variability and Outlier Detection,
2020--2025}
\author{Iowa Department of Health and Human Services}
\date{}
\begin{document}
\maketitle

\renewcommand*\contentsname{Table of contents}
{
\setcounter{tocdepth}{3}
\tableofcontents
}

```\{julia setup, include=FALSE\}

\#\textbar{} label: setup

\section{prepare to load packages}\label{prepare-to-load-packages}

using Pkg

\#= Ensure required packages are available Pkg.activate(``.'')
Pkg.instantiate()

only need to install packages the first time Pkg.add({[}``Tidier'',
``TidierPlots'', ``TidierDates'', ``DotEnv'', ``CSV'', ``DataFrames'',
``Quarto''{]}) =\#

\section{Load packages}\label{load-packages}

\section{Load packages}\label{load-packages-1}

using Tidier using TidierDates using Dates using TidierPlots using
DotEnv using CSV using DataFrames using Quarto

\section{Create .env file if it does not
exist}\label{create-.env-file-if-it-does-not-exist}

if !isfile(``.env'') write(``.env'', ``\,``\,''
IOWA\_TRAUMA\_REGISTRY\_COUNT\_PATH\_2020=
IOWA\_TRAUMA\_REGISTRY\_COUNT\_PATH\_2021=
IOWA\_TRAUMA\_REGISTRY\_COUNT\_PATH\_2022=
IOWA\_TRAUMA\_REGISTRY\_COUNT\_PATH\_2023=
IOWA\_TRAUMA\_REGISTRY\_COUNT\_PATH\_2024=
IOWA\_TRAUMA\_REGISTRY\_COUNT\_PATH\_2025=
IOWA\_TRAUMA\_REGISTRY\_COUNT\_PATH\_2026= ``\,``\,``) end

\section{Load .env file into ENV{[}{]}}\label{load-.env-file-into-env}

DotEnv.load!()

\section{Assign environment variables (no paths
yet)}\label{assign-environment-variables-no-paths-yet}

iowa\_trauma\_registry\_counts\_path\_2020 =
ENV{[}``IOWA\_TRAUMA\_REGISTRY\_COUNTS\_PATH\_2020''{]}
iowa\_trauma\_registry\_counts\_path\_2021 =
ENV{[}``IOWA\_TRAUMA\_REGISTRY\_COUNTS\_PATH\_2021''{]}
iowa\_trauma\_registry\_counts\_path\_2022 =
ENV{[}``IOWA\_TRAUMA\_REGISTRY\_COUNTS\_PATH\_2022''{]}
iowa\_trauma\_registry\_counts\_path\_2023 =
ENV{[}``IOWA\_TRAUMA\_REGISTRY\_COUNTS\_PATH\_2023''{]}
iowa\_trauma\_registry\_counts\_path\_2024 =
ENV{[}``IOWA\_TRAUMA\_REGISTRY\_COUNTS\_PATH\_2024''{]}
iowa\_trauma\_registry\_counts\_path\_2025 =
ENV{[}``IOWA\_TRAUMA\_REGISTRY\_COUNTS\_PATH\_2025''{]}
iowa\_trauma\_registry\_counts\_path\_2026 =
ENV{[}``IOWA\_TRAUMA\_REGISTRY\_COUNTS\_PATH\_2026''{]}

\begin{verbatim}

<!-- Load data -->
```{julia data_load, include=FALSE}

#| label: data_load

# Simple validation for missing paths
required_paths = [
    ("2020", iowa_trauma_registry_counts_path_2020),
    ("2021", iowa_trauma_registry_counts_path_2021),
    ("2022", iowa_trauma_registry_counts_path_2022),
    ("2023", iowa_trauma_registry_counts_path_2023),
    ("2024", iowa_trauma_registry_counts_path_2024),
    ("2025", iowa_trauma_registry_counts_path_2025),
    ("2026", iowa_trauma_registry_counts_path_2026)
]

for (yr, p) in required_paths
    if isempty(p)
        @warn "Path for $yr is empty. Update .env before attempting to load data."
    end
end

# Load all project data for anomaly detection via environment variable paths
 trauma_registry_counts_2020 = CSV.read(iowa_trauma_registry_counts_path_2020, DataFrame)
 trauma_registry_counts_2021 = CSV.read(iowa_trauma_registry_counts_path_2021, DataFrame)
 trauma_registry_counts_2022 = CSV.read(iowa_trauma_registry_counts_path_2022, DataFrame)
 trauma_registry_counts_2023 = CSV.read(iowa_trauma_registry_counts_path_2023, DataFrame)
 trauma_registry_counts_2024 = CSV.read(iowa_trauma_registry_counts_path_2024, DataFrame)
 trauma_registry_counts_2025 = CSV.read(iowa_trauma_registry_counts_path_2025, DataFrame)
 
\end{verbatim}

```\{julia data\_manipulation, include=FALSE\}

\#\textbar{} label: data\_manipulation

\section{union all rows from files and select the needed
variables}\label{union-all-rows-from-files-and-select-the-needed-variables}

trauma\_registry\_counts\_2020\_2025 = @chain
trauma\_registry\_counts\_2020 begin @bind\_rows(
trauma\_registry\_counts\_2021, trauma\_registry\_counts\_2022,
trauma\_registry\_counts\_2023, trauma\_registry\_counts\_2024,
trauma\_registry\_counts\_2025 ) @select(year, facility, total) end

\section{pivot columns wider}\label{pivot-columns-wider}

trauma\_registry\_counts\_2020\_2025\_pivot = @chain
trauma\_registry\_counts\_2020\_2025 begin @pivot\_wider( names\_from =
year, values\_from = total ) @mutate(
across(\texttt{2020}:\texttt{2025}, x -\textgreater{} coalesce.(x, 0)) )
@mutate( total\_2020\_2025 = \texttt{2020\_function} +
\texttt{2021\_function} + \texttt{2022\_function} +
\texttt{2023\_function} + \texttt{2024\_function} +
\texttt{2025\_function} ) @select(facility,
\texttt{2020\_function}:\texttt{2025\_function}, total\_2020\_2025)
@rename( \texttt{2020} = \texttt{2020\_function}, \texttt{2021} =
\texttt{2021\_function}, \texttt{2022} = \texttt{2022\_function},
\texttt{2023} = \texttt{2023\_function}, \texttt{2024} =
\texttt{2024\_function}, \texttt{2025} = \texttt{2025\_function} ) end

\section{finalize data by adding in differences between
years}\label{finalize-data-by-adding-in-differences-between-years}

trauma\_registry\_counts\_2020\_2025\_final = @chain
trauma\_registry\_counts\_2020\_2025\_pivot begin @mutate( facility =
replace.(facility, ``\x96'' =\textgreater{} ``--''), \# en dash,
diff\_2021 = \texttt{2021} - \texttt{2020}, diff\_2022 = \texttt{2022} -
\texttt{2021}, diff\_2023 = \texttt{2023} - \texttt{2022}, diff\_2024 =
\texttt{2024} - \texttt{2023}, diff\_2025 = \texttt{2025} -
\texttt{2024}, ) @mutate( mean\_records = round(mean.(c(\texttt{2020},
\texttt{2021}, \texttt{2022}, \texttt{2023}, \texttt{2024},
\texttt{2025})), digits = 3), var\_records = round(var.(c(\texttt{2020},
\texttt{2021}, \texttt{2022}, \texttt{2023}, \texttt{2024},
\texttt{2025})), digits = 3), sd\_records = round(std.(c(\texttt{2020},
\texttt{2021}, \texttt{2022}, \texttt{2023}, \texttt{2024},
\texttt{2025})), digits = 3), mean\_diff = round(mean.(c(diff\_2021,
diff\_2022, diff\_2023, diff\_2024, diff\_2025)), digits = 3), var\_diff
= round(var.(c(diff\_2021, diff\_2022, diff\_2023, diff\_2024,
diff\_2025)), digits = 3), sd\_diff = round(std.(c(diff\_2021,
diff\_2022, diff\_2023, diff\_2024, diff\_2025)), digits = 3) ) @mutate(
z\_score\_diff\_2021 = div(diff\_2021, sd\_diff), z\_score\_diff\_2022 =
div(diff\_2022, sd\_diff), z\_score\_diff\_2023 = div(diff\_2023,
sd\_diff), z\_score\_diff\_2024 = div(diff\_2024, sd\_diff),
z\_score\_diff\_2025 = div(diff\_2025, sd\_diff) ) @mutate(
anomaly\_2021 = abs(z\_score\_diff\_2021) \textgreater{} 1,
anomaly\_2022 = abs(z\_score\_diff\_2022) \textgreater{} 1,
anomaly\_2023 = abs(z\_score\_diff\_2023) \textgreater{} 1,
anomaly\_2024 = abs(z\_score\_diff\_2024) \textgreater{} 1,
anomaly\_2025 = abs(z\_score\_diff\_2025) \textgreater{} 1 ) @mutate(
any\_anomaly = any(c( anomaly\_2021, anomaly\_2022, anomaly\_2023,
anomaly\_2024, anomaly\_2025 )) ) end

\section{Before additional data manipulation and modeling, plot
differences}\label{before-additional-data-manipulation-and-modeling-plot-differences}

diff\_long = @chain trauma\_registry\_counts\_2020\_2025\_final begin
@select(facility, diff\_2021, diff\_2022, diff\_2023, diff\_2024,
diff\_2025) @pivot\_longer( diff\_2021:diff\_2025, names\_to = :year,
values\_to = :diff ) @mutate( year = replace.(year, ``diff\_''
=\textgreater{} ``\,``), sign = ifelse(diff .\textless{} 0,''negative'',
``positive'') ) end

\begin{verbatim}

# Introduction
Accurate and complete reporting to the Iowa Trauma Registry is essential for statewide injury surveillance, system evaluation, and trauma program performance improvement. Annual case counts submitted by each facility support a wide range of analytic and operational needs, including burden estimation, regional resource planning, and assessment of trauma care access. Substantial year‑to‑year variation in reporting can indicate changes in patient volume, modifications in facility workflows, electronic health record transitions, or potential data quality issues.

This report examines annual reporting differences from 2020 through 2025 for all participating Iowa trauma facilities. The analysis applies an anomaly‑detection method focused on  to identify facilities with reporting patterns that deviate from expected values. These methods include standardized difference scores, robust dispersion measures, and model‑based trend residuals. Facilities with large positive or negative deviations are flagged for further review to support quality assurance activities.

The objective is not to infer causation but to support registry data validation, highlight facilities requiring follow‑up, and strengthen the integrity of statewide trauma surveillance.

# Test table of 2025 results
```{julia test_table, echo=TRUE}

# subset the table with columns we want to see and fit
quarto_table = @chain trauma_registry_counts_2020_2025_final begin
    @select :facility, :mean_records, :sd_records, :mean_diff, :sd_diff, contains("2025")
end

PrettyTables.pretty_table(
    quarto_table,
    column_labels=names(quarto_table),
    show_row_number_column=false,
    title="Facilities with Anomalous "
)
\end{verbatim}




\end{document}
